\section{What a source-free bearing-diagnosis result should carry}
\label{sec:recommendations}

The measurements above suggest six pieces of evidence that separate a transferable result from a leaky one. None requires
new hardware or a larger dataset; together they cost a few additional runs.

\begin{enumerate}
  \item \textbf{Bearing-disjoint targets.} Report accuracy on physical bearings absent from source training. If the dataset
        cannot support this --- CWRU's healthy class is effectively one specimen, for example --- say so, and treat the
        published number as an upper bound rather than an estimate.
  \item \textbf{A distribution over splits, not one split.} Which bearings fall in the source set moved clean accuracy
        between \SI{36}{} and \SI{74}{\percent} here. A single split cannot distinguish a method effect from a draw effect;
        ten pre-registered draws are affordable and make the uncertainty visible.
  \item \textbf{Paired differences for any add-on module.} Gates, regularisers and losses should be compared against the
        base method trained from the \emph{same} source checkpoint with the same seed, and reported as the distribution of
        per-task differences. The relevant variance is that of the difference, which here is roughly half the variance of
        the accuracies themselves.
  \item \textbf{An oracle pseudo-label ceiling.} One extra run per configuration, replacing the pseudo-label rule by the
        true label, converts a bare gain into a share of what was achievable and exposes the part of the gap that lives in
        the source representation.
  \item \textbf{A non-adaptive shallow baseline under the identical split.} Cheap, and it calibrates the claim: on our
        splits it beat the adapted method.
  \item \textbf{False acceptance for any physics module, measured on healthy recordings from the target rig.} A gate that
        certifies faults on healthy machinery will inject confident wrong labels; this is measurable before adaptation and
        predicted the sign of the effect in every case we tested.
\end{enumerate}

For deployment the same list reads as a risk assessment. A diagnosis system that has only been validated with the monitored
bearing present in training has not been validated for the bearing it will meet in service; the drop measured here, on two
rigs and with three model families, was 22--41 points. Where a kinematic gate is available and its false acceptance has been
characterised on healthy machinery, it is a low-risk addition --- consistent, small and essentially never harmful in our
runs --- but it is not a substitute for source data that spans enough physical specimens. For supervised models, Vieira et al.\ \cite{vieira2026} report that the number
of unique training bearings is the decisive factor for robustness. Our splits hold that number fixed at three bearings per
class, so we cannot test it here; whether a larger source pool would close the gap measured in this paper is the obvious
next experiment, and our results are consistent with their finding without being evidence for it.
